\section{Introduction}

Federated learning allows multiple clients to collaboratively learn predictive models while keeping private training data local. In realistic systems, clients may differ not only in data distribution but also in model architecture, motivating \textbf{model-heterogeneous federated learning (HFL)}. Knowledge-distillation approaches such as FedDF communicate predictions rather than assuming parameter compatibility \citep{lin2020feddf}. Robust HFL further considers noisy or corrupted clients: RHFL studies noisy labels under model heterogeneity \citep{fang2022rhfl}, while AugHFL and its TPAMI extension RAHFL study model-heterogeneous collaboration under common data corruption \citep{fang2023aughfl,fang2025rahfl}. These works establish corruption as an important HFL robustness problem, but common corruption benchmarks and robust-learning pipelines typically evaluate corruption as a label-preserving nuisance rather than as a class-dependent training cue \citep{hendrycks2019corruptions,cubuk2020augmix}.

We study a more specific failure mode. Suppose one client observes most images of one class with blur and most images of another class with noise, while another client has a different class--corruption association. Corruption can then become predictive of the task label instead of acting only as nuisance variation. We call this controlled setting \textbf{class--corruption entanglement in HFL (CLE-HFL)}. The scientific question is not merely whether corruption lowers accuracy, but whether changing only the corruption operator moves predictions \emph{toward the classes associated with that operator during training}. Corrupted accuracy alone cannot distinguish this directional shortcut from ordinary image degradation, label imbalance, model capacity limits, or non-IID optimization.

This motivates a target-aligned diagnostic. Within-view consistency, including AugMix-style Jensen--Shannon consistency, can make multiple augmentations around the same corrupted observation agree while leaving cross-operator class routing intact \citep{cubuk2020augmix}. Likewise, a pseudo-environment statistic can improve without identifying whether the model has reduced reliance on the true class--corruption association. We therefore evaluate the same semantic source under multiple corruption operators and define \textbf{Directional Shortcut Alignment (DSA)} as a probability-mass contrast toward the classes pre-registered as bound to the applied operator. DSA is intentionally narrower than a general causal representation claim: under semantic preservation, a valid operator intervention, pre-registered binding, source-level pairing, and strict evaluation isolation, it identifies a binding-direction operator-response contrast. We further show that proxy-only improvements cannot generally certify reduced true CLE dependence when the proxy error channel is unconstrained, making a target-aligned behavioral endpoint necessary under our evidence protocol.

Using DSA, we next ask where the shortcut forms. In a seed-0 matched HFL-versus-Local factorial, the pooled strong-CLE effect is close between HFL and Local training, while the additional communication contrast is much smaller. We therefore describe the observed formation pattern as \textbf{predominantly local-first} in this controlled setting and use that result to motivate intervention at the client objective rather than redesigning communication. The intervention combines a \textbf{coarse family-level Public Environment Witness (PEW)}, trained only on public synthetic corruptions, with \textbf{Balanced Environment Risk (BER)}, which reduces the effective risk-mass advantage of large pseudo-environment supports within each task class. PEW is explicitly \textbf{taxonomy-assisted}, not taxonomy-free, and the training method never reads private true corruption operators or families.

We evaluate the resulting evidence chain rather than presenting PEW or BER as isolated architectural novelties. A held-out 40-round comparison uses matched ERM, a protocol-matched CVaR-DRO tail-risk control, PEW+GroupDRO with the same PEW grouping information, and PEW+BER. Additional controlled experiments change only the class--operator binding map or the HFL communication base. The completed evidence supports strong shortcut suppression, but not a universal no-harm claim: the FedDF-fidelity utility gate fails, and CVaR-DRO retains lower DSA for some individual clients on the held-out map. Our contributions are therefore fourfold:

\begin{enumerate}
\item \textbf{CLE-HFL problem formulation.} We formulate class--corruption entanglement in model-heterogeneous federated learning, where client-specific class--corruption associations can turn corruption operators into directional class cues.
\item \textbf{Paired diagnosis and identification boundary.} We introduce source-paired operator counterfactual evaluation and \textbf{Directional Shortcut Alignment (DSA)}, characterize its identification assumptions and source-level inference, and establish a proxy non-identifiability result showing why generic proxy improvements cannot certify reduced CLE harm without additional assumptions.
\item \textbf{Local-versus-communication formation analysis.} Using a seed-0 matched HFL-versus-Local factorial, we find that the observed directional shortcut is predominantly local-first in the controlled setting, with communication contributing a smaller pooled additional effect.
\item \textbf{Taxonomy-assisted local mitigation and controlled validation.} Guided by this finding, we develop a coarse family-level taxonomy-assisted local intervention combining PEW and BER, and evaluate it against matched ERM, tail-risk, and group-robustness controls across controlled binding maps and HFL communication bases while explicitly reporting utility and taxonomy boundaries.
\end{enumerate}

We do not claim to be the first to study spurious correlation in federated learning, environment inference, group reweighting, or counterfactual shortcut evaluation \citep{wang2024pflspurious,tang2024fedpin,ma2024fedcd,creager2021eiil,sagawa2020groupdro,vigneshwaran2026counterfactual}. Our focus is the joint instantiation of a client-specific class--corruption directional failure mode, a binding-aligned paired diagnostic, local-versus-communication attribution, and a structure-matched local intervention in model-heterogeneous FL.
